\section{The Jacobi Equation}

Let $I,K\subseteq\mathbb R$ be intervals, let $\gamma:I\to M$ be a geodesic,
and let $\Gamma:K\times I\to M$ be a smooth variation.  It is a
\emph{variation through geodesics} if every main curve
$\Gamma_s(t)=\Gamma(s,t)$ is a geodesic.

\begin{theorem}[The Jacobi Equation]
\label{thm:jacobi-equation}
\dcref{ch10:10.1}
Let $(M,g)$ be Riemannian or pseudo-Riemannian, let $\gamma$ be a geodesic,
and let $J$ be the variation field of a variation through geodesics.  Then
\[
D_t^2J+R(J,\gamma')\gamma'=0. \quad (10.1)
\]
A vector field along $\gamma$ satisfying this equation is a \emph{Jacobi
field}.
\end{theorem}

\begin{proof}
\sketch
Set $T=\partial_t\Gamma$ and $S=\partial_s\Gamma$.  Since each main curve is
geodesic, $D_tT=0$.  Commuting covariant derivatives and using the symmetry
lemma gives
\[
0=D_sD_tT=D_tD_sT+R(S,T)T=D_t^2S+R(S,T)T.
\]
At $s=0$, this is the stated equation for $J=S(0,\mathord\cdot)$.
\end{proof}

\begin{proposition}[Existence and Uniqueness of Jacobi Fields]
\label{prop:existence-uniqueness-jacobi-fields}
\dcref{ch10:10.2}
\uses{thm:linear-ode-existence-uniqueness}
Let $I\subseteq\mathbb R$ be an interval, let $\gamma:I\to M$ be a geodesic,
let $a\in I$, and put $p=\gamma(a)$.  For every $v,w\in T_pM$, there is a
unique Jacobi field along $\gamma$ satisfying
\[
J(a)=v,\qquad D_tJ(a)=w.
\]
\end{proposition}

\begin{proof}
\sketch
Choose a parallel frame $(E_i)$ along $\gamma$, write
$\gamma'=y^iE_i$ and $J=J^iE_i$, and set $W^i=\dot J^i$.  The Jacobi equation
becomes the first-order linear system
\[
\dot J^i=W^i,
\qquad
\dot W^i=-R_{jki}^{\ell}(\gamma(t))J^\ell y^ky^j.
\]
The linear ODE existence--uniqueness theorem gives a unique smooth solution on
all of $I$ for $J^i(a)=v^i$ and $W^i(a)=w^i$, which is the required field.
\end{proof}

For a geodesic $\gamma$, write $\mathfrak J(\gamma)$ for the vector space of
Jacobi fields along it.

\begin{corollary}
\label{cor:jacobi-space-dimension}
\dcref{ch10:10.3}
\uses{prop:existence-uniqueness-jacobi-fields}
If $(M,g)$ has dimension $n$ and $\gamma$ is any geodesic, then
$\mathfrak J(\gamma)$ is a $2n$-dimensional linear space.
\end{corollary}

\begin{proof}
\sketch
The Jacobi equation is linear.  Evaluation at any $a\in I$ defines
\[
\mathfrak J(\gamma)\longrightarrow T_{\gamma(a)}M\oplus T_{\gamma(a)}M,
\qquad
J\longmapsto(J(a),D_tJ(a)),
\]
which is an isomorphism by the preceding proposition.
\end{proof}

\begin{proposition}[Jacobi Fields as Variation Fields]
\label{prop:jacobi-field-variation-field}
\dcref{ch10:10.4}
\uses{thm:jacobi-equation}
Let $(M,g)$ be Riemannian or pseudo-Riemannian and let $\gamma:I\to M$ be a
geodesic.  If $M$ is complete or $I$ is compact, every Jacobi field along
$\gamma$ is the variation field of a variation through geodesics.
\end{proposition}

\begin{proof}
\sketch
Translate the parameter so that $0\in I$, set $p=\gamma(0)$ and
$v=\gamma'(0)$, and choose a curve $\sigma$ and vector field $V$ along it with
\[
\sigma(0)=p,\quad \sigma'(0)=J(0),\quad
V(0)=v,\quad D_sV(0)=D_tJ(0).
\]
Define
\begin{equation}
\Gamma(s,t)=\exp_{\sigma(s)}(tV(s)). \tag{10.2}
\end{equation}
Completeness, or compactness of $I$ and openness of the exponential domain,
makes this defined on a product neighborhood of $\{0\}\times I$.  It obeys
\begin{align}
\Gamma(0,t)&=\gamma(t), \tag{10.3}\\
\Gamma(s,0)&=\sigma(s). \tag{10.4}
\end{align}
Its variation field $W$ is Jacobi, $W(0)=J(0)$, and the symmetry lemma gives
\[
D_tW(0)=D_s\partial_t\Gamma(0,0)=D_sV(0)=D_tJ(0).
\]
Jacobi uniqueness yields $W=J$.
\end{proof}

\begin{proposition}[Local Isometry Invariance of Jacobi Fields]
\label{prop:local-isometry-invariance-jacobi-fields}
\dcref{ch10:10.5}
If $\varphi:(M,g)\to(\widetilde M,\widetilde g)$ is a local isometry,
$\widetilde\gamma=\varphi\circ\gamma$, and
$d\varphi_{\gamma(t)}J(t)=\widetilde J(t)$, then $J$ is Jacobi if and only if
$\widetilde J$ is Jacobi.
\end{proposition}

\begin{exercise}
\label{exer:prove-local-isometry-invariance}
\dcref{ch10:10.6}
\uses{prop:local-isometry-invariance-jacobi-fields}
Prove the preceding proposition.
\end{exercise}

\section{Basic Computations with Jacobi Fields}

\subsection{Tangential and Normal Jacobi Fields}

Along every geodesic, $J_0(t)=\gamma'(t)$ and $J_1(t)=t\gamma'(t)$ are Jacobi
fields.  They arise from the reparametrization variations
$\gamma(s+t)$ and $\gamma((1+s)t)$ whenever defined.

For a regular curve, let
$T^\top_{\gamma(t)}M=\operatorname{span}(\gamma'(t))$ and let
$T^\perp_{\gamma(t)}M$ be its orthogonal complement.  A field is
\emph{tangential} or \emph{normal} when its values lie in these subspaces,
respectively.  Write $\mathfrak X^\top(\gamma)$ and
$\mathfrak X^\perp(\gamma)$ for the corresponding spaces of fields, and
$\mathfrak J^\top(\gamma)$ and $\mathfrak J^\perp(\gamma)$ for the Jacobi
subspaces.

\begin{proposition}[Equivalent Conditions for a Normal Jacobi Field]
\label{prop:normal-jacobi-equivalent-conditions}
\dcref{ch10:10.7}
Let $(M,g)$ be Riemannian or pseudo-Riemannian, let $\gamma:I\to M$ be a
geodesic, and let $J$ be a Jacobi field.  The following are equivalent:
\begin{enumerate}
\item[(a)] $J$ is normal.
\item[(b)] $J$ is orthogonal to $\gamma'$ at two distinct points.
\item[(c)] $J$ and $D_tJ$ are orthogonal to $\gamma'$ at one point.
\item[(d)] $J$ and $D_tJ$ are orthogonal to $\gamma'$ everywhere.
\end{enumerate}
\end{proposition}

\begin{proof}
\sketch
Let $f(t)=\langle J(t),\gamma'(t)\rangle$.  Since $D_t\gamma'=0$, the Jacobi
equation and curvature symmetries give
\[
f''=\langle D_t^2J,\gamma'\rangle
=-\mathrm{Rm}(J,\gamma',\gamma',\gamma')=0.
\]
Thus $f$ is affine, while
$f'=\langle D_tJ,\gamma'\rangle$.  An affine function vanishes at two points
if and only if it vanishes identically, equivalently if its value and
derivative vanish at one point.  This is exactly (a)--(d).
\end{proof}

\begin{corollary}[Tangential--Normal Jacobi Decomposition]
\label{cor:jacobi-field-decomposition}
\dcref{ch10:10.8}
\uses{cor:jacobi-space-dimension,prop:normal-jacobi-equivalent-conditions}
Suppose $(M,g)$ is Riemannian or pseudo-Riemannian and $\gamma$ is a
nonconstant non-null geodesic.  Then
$\mathfrak J^\perp(\gamma)$ has dimension $2n-2$,
$\mathfrak J^\top(\gamma)$ has dimension $2$, and
\[
\mathfrak J(\gamma)=\mathfrak J^\perp(\gamma)\oplus\mathfrak J^\top(\gamma).
\]
Every Jacobi field has a unique tangential-plus-normal decomposition.
\end{corollary}

\begin{proof}
\sketch
Initial-data evaluation identifies $\mathfrak J(\gamma)$ with two copies of
$T_{\gamma(a)}M$.  Proposition 10.7 identifies the normal subspace with pairs
orthogonal to $\gamma'(a)$, giving dimension $2n-2$.  The fields $J_0,J_1$
span a two-dimensional tangential subspace.  For a non-null geodesic, the
tangential and normal subspaces intersect trivially, so their dimensions sum
to $2n$ and give the direct sum.
\end{proof}

\section{Jacobi Fields Vanishing at a Point}

\begin{lemma}[Variation Formula for a Jacobi Field Vanishing at the Origin]
\label{lem:jacobi-field-variation}
\dcref{ch10:10.9}
\uses{prop:jacobi-field-variation-field}
Let $I$ contain $0$, let $\gamma:I\to M$ be a geodesic, and let $J$ be Jacobi
with $J(0)=0$.  If $M$ is complete or $I$ is compact, then with
$p=\gamma(0)$, $v=\gamma'(0)$, and $w=D_tJ(0)$,
\begin{equation}
\Gamma(s,t)=\exp_p(t(v+sw)) \quad (10.5)
\end{equation}
is a variation through geodesics having variation field $J$.
\end{lemma}

\begin{proof}
\sketch
In the construction of Proposition 10.4, take the base curve
$\sigma(s)\equiv p$ and the vector field $V(s)=v+sw\in T_pM$.  The resulting
variation is (10.5), and its field has the same initial value and derivative
as $J$; uniqueness identifies them.
\end{proof}

\begin{proposition}[Jacobi Fields Vanishing at a Point]
\label{prop:jacobi-fields-vanishing-at-point}
\dcref{ch10:10.10}
Let $(M,g)$ be Riemannian or pseudo-Riemannian, let $p\in M$, and let
$\gamma:I\to M$ satisfy $0\in I$ and $\gamma(0)=p$.  If
$v=\gamma'(0)$ and $J(0)=0$, $D_tJ(0)=w$, then
\[
J(t)=d(\exp_p)_{tv}(tw), \quad (10.6)
\]
under $T_{tv}(T_pM)\cong T_pM$.  In normal coordinates on a neighborhood
containing $\gamma(I)$,
\begin{equation}
J(t)=tw^i\partial_i\bigg|_{\gamma(t)}, \tag{10.7}
\end{equation}
where $w=w^i\partial_i|_0$.
\end{proposition}

\begin{proof}
\sketch
Restrict to any compact interval containing $0$ and apply Lemma 10.9.  The
chain rule applied to (10.5) gives (10.6); exhaustion by compact intervals
extends it across $I$.  In normal coordinates,
\[
\Gamma(s,t)=\bigl(t(v^1+sw^1),\ldots,t(v^n+sw^n)\bigr),
\]
so differentiating in $s$ gives (10.7).
\end{proof}

\begin{corollary}
\label{cor:jacobi-field-tangent-vectors}
\dcref{ch10:10.11}
\uses{prop:jacobi-fields-vanishing-at-point}
If $U$ is a normal neighborhood of $p$, every vector in $T_qM$ for
$q\in U\setminus\{p\}$ is the value at $q$ of a Jacobi field along a radial
geodesic that vanishes at $p$.
\end{corollary}

\begin{proof}
\sketch
In normal coordinates write $q=(q^1,\ldots,q^n)$ and
$w=w^i\partial_i|_q$.  The radial geodesic
$\gamma(t)=(tq^1,\ldots,tq^n)$ joins $p$ to $q$, and (10.7) gives
$J(t)=tw^i\partial_i|_{\gamma(t)}$, with $J(0)=0$ and $J(1)=w$.
\end{proof}

\section{Jacobi Fields in Constant-Curvature Spaces}

For $c\in\mathbb R$, define
\begin{equation}
s_c(t)=\begin{cases}
t,&c=0,\\
R\sin(t/R),&c=1/R^2>0,\\
R\sinh(t/R),&c=-1/R^2<0.
\end{cases} \tag{10.8}
\end{equation}

\begin{proposition}[Jacobi Fields in Constant Curvature]
\label{prop:jacobi-constant-curvature}
\dcref{ch10:10.12}
Suppose $(M,g)$ has constant sectional curvature $c$ and $\gamma$ is
unit-speed.  The normal Jacobi fields along $\gamma$ with $J(0)=0$ are
\begin{equation}
J(t)=ks_c(t)E(t), \tag{10.9}
\end{equation}
where $E$ is any parallel unit normal field and $k$ is constant.  Moreover,
\begin{equation}
D_tJ(0)=kE(0), \tag{10.10}
\end{equation}
and
\begin{equation}
|J(t)|=|s_c(t)|\,|D_tJ(0)|. \tag{10.11}
\end{equation}
\end{proposition}

\begin{proof}
\sketch
The constant-curvature formula gives
$R(v,w)x=c(\langle w,x\rangle v-\langle v,x\rangle w)$.  For a normal
field, the Jacobi equation reduces to
\begin{equation}
\begin{aligned}
0&=D_t^2J+c(\langle\gamma',\gamma'\rangle J
-\langle J,\gamma'\rangle\gamma')\\
&=D_t^2J+cJ.
\end{aligned} \tag{10.12}
\end{equation}
Writing $J=uE$ with $E$ parallel and normal reduces this to
$u''+cu=0$.  The solutions with $u(0)=0$ are the multiples of $s_c$; these
give all normal fields by initial-data dimension.  Since $s_c'(0)=1$,
(10.10) holds, and unit length of $E$ gives (10.11).
\end{proof}

Let $\pi:\mathbb R^n\setminus\{0\}\to S^{n-1}$ be radial projection,
\begin{equation}
\pi(x)=\frac{x}{|x|}. \tag{10.13}
\end{equation}
Let $\check g$ be the unit round metric on $S^{n-1}$ and define
\begin{equation}
\hat g=\pi^*\check g. \quad (10.14)
\end{equation}

\begin{lemma}[Polar Form of the Euclidean Metric]
\label{lem:metric-relation-euclidean}
\dcref{ch10:10.13}
\uses{ex:warped-products,prob:warped-product-isometry}
On $\mathbb R^n\setminus\{0\}$, the Euclidean metric $\bar g$ satisfies
\begin{equation}
\bar g=dr^2+r^2\hat g. \quad (10.15)
\end{equation}
where $r(x)=|x|$.
\end{lemma}

\begin{proof}
\sketch
The map
\begin{equation}
\Phi(\rho,\omega)=\rho\omega \quad (10.16)
\end{equation}
is an isometry from the warped product
$\mathbb R^+\times_\rho S^{n-1}$ with metric
$d\rho^2\oplus\rho^2\check g$ to Euclidean space minus the origin.  Since
$\Phi^{-1}(x)=(r(x),\pi(x))$, pulling back this identity gives (10.15).
\end{proof}

\begin{theorem}[Constant-Curvature Metrics in Normal Coordinates]
\label{thm:constant-curvature-normal-coords}
\dcref{ch10:10.14}
\uses{thm:minimizing-geodesic-unit-speed,cor:jacobi-field-decomposition}
Suppose $(M,g)$ has constant sectional curvature $c$.  In normal coordinates
around $p$, let $r$ be the radial distance and let $\hat g$ be (10.14).  On
$U\setminus\{p\}$,
\begin{equation}
g=dr^2+s_c(r)^2\hat g. \quad (10.17)
\end{equation}
\end{theorem}

\begin{proof}
\sketch
Let $g_c$ denote the right side of (10.17), and let $\bar g$ be the Euclidean
metric in the normal coordinates.  All three metrics make $\partial_r$ unit
and orthogonal to the $r$-level sets, so compare tangent vectors to one such
level set.  If $q$ lies on $r=b$ and $w$ is tangent there, the radial geodesic
$\gamma(t)=(tq^1/b,\ldots,tq^n/b)$ and the Jacobi field
\begin{equation}
J(t)=\frac{t}{b}w^i\partial_i\bigg|_{\gamma(t)} \quad (10.18)
\end{equation}
have $J(b)=w$ and $D_tJ(0)=b^{-1}w^i\partial_i|_p$.  They are normal, so
Proposition 10.12 gives
\[
|w|_g^2=s_c(b)^2|D_tJ(0)|_g^2
=\frac{s_c(b)^2}{b^2}|w_p|_g^2
=\frac{s_c(b)^2}{b^2}|w|_{\bar g}^2
=|w|_{g_c}^2,
\]
where $w_p=w^i\partial_i|_p$ and the middle equality uses that normal
coordinates make $g_p=\bar g_p$.  Thus (10.15) and polarization prove
$g=g_c$.
\end{proof}

\begin{corollary}[Local Uniqueness of Constant-Curvature Metrics]
\label{cor:local-uniqueness-constant-curvature}
\dcref{ch10:10.15}
\uses{thm:constant-curvature-normal-coords}
If $(M,g)$ and $(\widetilde M,\widetilde g)$ have the same dimension and
constant sectional curvature $c$, then every $p\in M$ and
$\widetilde p\in\widetilde M$ have neighborhoods related by an isometry.
\end{corollary}

\begin{proof}
\sketch
Choose equal-radius small geodesic balls around $p$ and $\widetilde p$.  By
(10.17), their normal-coordinate metric expressions agree, and the coordinate
identification of the Euclidean balls extends across the origins because both
metrics equal $\delta_{ij}$ there.
\end{proof}

\begin{corollary}[Constant-Curvature Metrics as Warped Products]
\label{cor:constant-curvature-warped-products}
\dcref{ch10:10.16}
Suppose $U$ is a geodesic ball of radius $b$ in a constant-curvature
Riemannian manifold.  Then $U\setminus\{p\}$ is isometric to
$(0,b)\times_{s_c}S^{n-1}$, where $S^{n-1}$ has the unit round metric
$\check g$.
\end{corollary}

\begin{proof}
\sketch
Use (10.17) in normal coordinates and the map (10.16).  Since
$\pi\circ\Phi$ restricts to the identity on each sphere
$\{\rho\}\times S^{n-1}$,
\[
\Phi^*g=d\rho^2\oplus s_c(\rho)^2\check g.
\]
\end{proof}

\begin{corollary}[Polar Decomposition of Integrals]
\label{cor:polar-decomposition-integrals}
\dcref{ch10:10.17}
\uses{cor:constant-curvature-warped-products}
Let $U$ be an open or closed geodesic ball of radius $b$ in a
constant-curvature $n$-manifold, and let $f:U\to\mathbb R$ be bounded and
integrable.  Then
\[
\int_U f\,dV_g
=\int_{S^{n-1}}\int_0^b f\circ\Phi(\rho,\omega)
s_c(\rho)^{n-1}\,d\rho\,dV_{\check g},
\]
where $\Phi$ is (10.16).
\end{corollary}

\begin{proof}
\sketch
The ball boundary and its center have measure zero, so use the warped-product
description on $(0,b)\times S^{n-1}$.  The volume density of
$d\rho^2\oplus s_c(\rho)^2\check g$ is
$s_c(\rho)^{n-1}d\rho\,dV_{\check g}$, giving the formula.
\end{proof}
\section{Locally Symmetric Spaces}

A Riemannian manifold is \emph{locally symmetric} if every point has a
neighborhood admitting a point reflection at that point.

\begin{lemma}
\label{lem:locally-symmetric-cartan}
\dcref{ch10:10.18}

Let $(M,g)$ have parallel curvature tensor. Suppose $p,\widetilde p\in M$ and
$A:T_pM\to T_{\widetilde p}M$ satisfy
\[
A^*g_{\widetilde p}=g_p,
\qquad
A^*Rm_{\widetilde p}=Rm_p.
\]
Then some neighborhood $U$ of $p$ admits a local isometry
$\varphi:U\to M$ with
$\varphi(p)=\widetilde p$ and $d\varphi_p=A$. If $M$ is complete, $U$ may be
any normal neighborhood of $p$.
\end{lemma}

\begin{proof}
\sketch If $M$ is complete, let $U$ be any normal neighborhood of $p$;
otherwise choose $\varepsilon>0$ so that both
$B_\varepsilon(p)$ and $B_\varepsilon(\widetilde p)$ are geodesic balls and
put $U=B_\varepsilon(p)$. Define
\[
\varphi=\exp_{\widetilde p}\circ A\circ\exp_p^{-1}.
\]
Then $\varphi(p)=\widetilde p$ and $d\varphi_p=A$.

For $q=\exp_p(v)\in U\setminus\{p\}$ and $x\in T_qM$, choose the Jacobi field
$J$ along $\gamma(t)=\exp_p(tv)$ with $J(0)=0$, $J(1)=x$, and write
$D_tJ(0)=w$. Along
$\widetilde\gamma(t)=\exp_{\widetilde p}(tA v)$, let
$\widetilde J$ be the Jacobi field with
$\widetilde J(0)=0$ and $D_t\widetilde J(0)=Aw$. The exponential-map
variation formula gives
\[
\widetilde J(1)=d\varphi_q(J(1))=d\varphi_q(x).
\]

Choose parallel orthonormal frames $(E_i)$ and $(\widetilde E_i)$ along the
two geodesics with $\widetilde E_i(0)=A(E_i(0))$. Parallelism of $Rm$ implies
parallelism of $R$ because covariant differentiation commutes with raising
indices. Hence the
coefficients
\[
R(E_i,E_j)E_k=R_{ijk}^{\,l}E_l
\]
are constant along $\gamma$, and the corresponding coefficients along
$\widetilde\gamma$ are the same constants by $A^*Rm_{\widetilde p}=Rm_p$.
Thus the component functions of $J$ and $\widetilde J$ solve the same Jacobi
ODE with initial data
\[
J^i(0)=\widetilde J^i(0)=0,
\qquad
(J^i)'(0)=(\widetilde J^i)'(0)=w^i.
\]
ODE uniqueness gives $J^i=\widetilde J^i$, so
$|d\varphi_q(x)|=|\widetilde J(1)|=|J(1)|=|x|$. Polarization shows that
$d\varphi_q$ preserves inner products, hence $\varphi$ is a local isometry.
\end{proof}

\begin{theorem}[Characterization of Locally Symmetric Spaces]
\label{thm:characterization-locally-symmetric-spaces}
\dcref{ch10:10.19}
\uses{prob:symmetric-space-parallel-curvature,lem:locally-symmetric-cartan}

A Riemannian manifold is locally symmetric if and only if its curvature
tensor is parallel.

\begin{proof}
\sketch Problem~7-3 proves that local symmetry implies $\nabla Rm=0$.
Conversely, fix $p$ and apply
Lemma~\ref{lem:locally-symmetric-cartan} with $\widetilde p=p$ and
$A=-\operatorname{Id}_{T_pM}$. Since $Rm_p$ is $4$-linear,
\[
A^*g_p=g_p,
\qquad
(A^*Rm_p)(w,x,y,z)=Rm_p(-w,-x,-y,-z)=Rm_p(w,x,y,z).
\]
The lemma gives a local isometry fixing $p$ with differential
$-\operatorname{Id}$. Shrinking to a sufficiently small geodesic ball makes
it the point reflection of that ball, proving local symmetry.
\end{proof}
\end{theorem}

\section{Conjugate Points}

Let $(M,g)$ be Riemannian or pseudo-Riemannian, let $\gamma:I\to M$ be a
geodesic, and let $p=\gamma(a)$ and $q=\gamma(b)$. The points $p,q$ are
\emph{conjugate along $\gamma$} if a nonzero Jacobi field along $\gamma$
vanishes at both $a$ and $b$. Their \emph{order}, or \emph{multiplicity}, is
the dimension of the space of such fields. This multiplicity is at most
$n-1$: the Jacobi fields vanishing at $a$ form an $n$-dimensional space, and
a nonzero tangential Jacobi field cannot vanish twice. The bound is attained
by antipodal points along a geodesic in $S^n(R)$, where the multiplicity is
$n-1$.

\begin{proposition}[Critical Points and Conjugacy]
\label{prop:critical-points-conjugate}
\dcref{ch10:10.20}
\lean{LeeLib.Ch10.criticalPoints_conjugate}
\leanok

Let $(M,g)$ be Riemannian or pseudo-Riemannian, let $p\in M$, and let
$v\in\mathcal E_p\subseteq T_pM$. Set
\[
\gamma(t)=\exp_p(tv)\quad(0\leq t\leq1),
\qquad
q=\gamma(1)=\exp_p(v).
\]
Then $v$ is a critical point of $\exp_p$ if and only if $q$ is conjugate to
$p$ along $\gamma$.

\begin{proof}
\sketch For $w\in T_pM$, use the geodesic variation
\[
\Gamma_w(s,t)=\exp_p\bigl(t(v+sw)\bigr)
\]
with variation field $J_w$. Its endpoint and initial-derivative data are
\[
J_w(0)=0,\qquad D_tJ_w(0)=w,\qquad
J_w(1)=d(\exp_p)_v(w).
\]
Every Jacobi field vanishing at $0$ is uniquely $J_w$ for
$w=D_tJ(0)$, and it is nonzero exactly when $w\neq0$. Thus
$\ker d(\exp_p)_v$ is nontrivial exactly when a nonzero Jacobi field vanishes
at both endpoints.
\end{proof}
\end{proposition}

The \emph{two-point boundary problem} along
$\gamma:[a,b]\to M$ asks, for prescribed
$v\in T_{\gamma(a)}M$ and $w\in T_{\gamma(b)}M$, for a Jacobi field $J$
satisfying $J(a)=v$ and $J(b)=w$.

\begin{proposition}[The Two-Point Boundary Problem for Jacobi Fields]
\label{prop:two-point-boundary-problem}
\dcref{ch10:10.21}
\uses{prop:existence-uniqueness-jacobi-fields,cor:jacobi-field-tangent-vectors,prop:critical-points-conjugate}

Let $(M,g)$ be Riemannian or pseudo-Riemannian and let
$\gamma:[a,b]\to M$ be a geodesic segment. The two-point boundary problem
has a unique solution for every
$v\in T_{\gamma(a)}M$ and $w\in T_{\gamma(b)}M$ if and only if
$\gamma(a)$ and $\gamma(b)$ are not conjugate along $\gamma$.
\end{proposition}

\section{The Second Variation Formula}

In this section $(M,g)$ is Riemannian. A variation is \emph{normal} if its
variation field is normal to the underlying curve.

\begin{theorem}[Second Variation Formula]
\label{thm:second-variation-formula}
\dcref{ch10:10.22}
\uses{thm:first-variation-formula,prop:curvature-commutator-failure}

Let $\gamma:[a,b]\to M$ be a unit-speed geodesic segment, let
$\Gamma:J\times[a,b]\to M$ be a proper variation, and let $V$ be its
variation field. Then
\begin{equation}
\left.\frac{d^2}{ds^2}\right|_{s=0}L_g(\Gamma_s)
=\int_a^b\left(
|D_tV^\perp|^2-Rm(V^\perp,\gamma',\gamma',V^\perp)
\right)\,dt,
\tag{10.19}
\end{equation}
where $V^\perp$ is the component of $V$ normal to $\gamma'$.

\begin{proof}
\sketch Put $T=\partial_t\Gamma$ and $S=\partial_s\Gamma$, and let
$(a_0,\ldots,a_k)$ be an admissible partition. On a smooth rectangle, twice
differentiating length, using $D_sT=D_tS$ and the curvature commutator, and
then setting $s=0$ gives
\begin{align}
\left.\frac{d^2}{ds^2}\right|_{s=0}
L_g(\Gamma_s|_{[a_{i-1},a_i]})
&=\left.\int_{a_{i-1}}^{a_i}
\bigl(
\langle D_tD_sS,T\rangle-Rm(S,T,T,S)\notag\\
&\qquad\qquad+|D_tS|^2-\langle D_tS,T\rangle^2
\bigr)\,dt\right|_{s=0}.
\tag{10.20}
\end{align}
Because $D_tT=D_t\gamma'=0$ at $s=0$,
\begin{equation}
\int_{a_{i-1}}^{a_i}\langle D_tD_sS,T\rangle\,dt
=\left.\langle D_sS,T\rangle\right|_{t=a_{i-1}}^{t=a_i}.
\tag{10.21}
\end{equation}
Properness gives $S=D_sS=0$ at $a,b$, and $D_sS$ is continuous across the
partition points, so these boundary terms cancel after summing over $i$.
Consequently,
\begin{align}
\left.\frac{d^2}{ds^2}\right|_{s=0}L_g(\Gamma_s)
&=\int_a^b\left(
|D_tV|^2-\langle D_tV,\gamma'\rangle^2
-Rm(V,\gamma',\gamma',V)
\right)\,dt.
\tag{10.22}
\end{align}

Decompose
\[
V=V^\top+V^\perp,\qquad
V^\top=\langle V,\gamma'\rangle\gamma'.
\]
Since $D_t\gamma'=0$,
\[
D_tV^\top=(D_tV)^\top,\qquad
D_tV^\perp=(D_tV)^\perp,
\]
and hence
\[
|D_tV|^2
=\langle D_tV,\gamma'\rangle^2+|D_tV^\perp|^2.
\]
Curvature antisymmetry gives
$Rm(V,\gamma',\gamma',V)
=Rm(V^\perp,\gamma',\gamma',V^\perp)$. Substitution into (10.22) proves
(10.19).
\end{proof}
\end{theorem}

For a geodesic segment $\gamma:[a,b]\to M$, its \emph{index form} is the
symmetric bilinear form on normal vector fields
\[
I(V,W)=\int_a^b
\left(
\langle D_tV,D_tW\rangle-Rm(V,\gamma',\gamma',W)
\right)\,dt.
\qquad (10.23)
\]

\begin{corollary}
\label{cor:second-variation-minimizing-geodesic}
\dcref{ch10:10.23}
\uses{thm:second-variation-formula}

Let $\gamma:[a,b]\to M$ be a unit-speed geodesic, let $\Gamma$ be a proper
normal variation, and let $V$ be its variation field. The second variation is
$I(V,V)$. If $\gamma$ is minimizing, then
$I(V,V)\geq0$ for every proper normal vector field $V$ along $\gamma$.
\end{corollary}

\begin{proposition}
\label{prop:index-form-jacobi-equation}
\dcref{ch10:10.24}

Let $\gamma:[a,b]\to M$ be a geodesic segment. For piecewise smooth normal
vector fields $V,W$ along $\gamma$,
\[
\begin{aligned}
I(V,W)
&=-\int_a^b
\left\langle D_t^2V+R(V,\gamma')\gamma',W\right\rangle\,dt\\
&\quad+\left.\langle D_tV,W\rangle\right|_{t=a}^{t=b}
-\sum_{i=1}^{k-1}\left\langle\Delta_iD_tV,W(a_i)\right\rangle,
\end{aligned}
\qquad (10.24)
\]
where $(a_0,\ldots,a_k)$ is an admissible partition for $V,W$, and
$\Delta_iD_tV$ is the jump of $D_tV$ at $a_i$.

\begin{proof}
\sketch On each smooth subinterval,
\[
\frac d{dt}\langle D_tV,W\rangle
=\langle D_t^2V,W\rangle+\langle D_tV,D_tW\rangle.
\]
Integrating gives
\[
\int_{a_{i-1}}^{a_i}\langle D_tV,D_tW\rangle\,dt
=-\int_{a_{i-1}}^{a_i}\langle D_t^2V,W\rangle\,dt
+\left.\langle D_tV,W\rangle\right|_{a_{i-1}}^{a_i}.
\]
Sum over $i$. Continuity of $W$ makes the interior endpoint contributions
equal the negative jump terms in (10.24), and adding the curvature term from
$I$ proves the formula.
\end{proof}
\end{proposition}

\begin{corollary}
\label{cor:jacobi-field-index-form}
\dcref{ch10:10.25}

Let $\gamma$ be a geodesic segment and let $V$ be a proper normal piecewise
smooth vector field along it. Then
$I(V,W)=0$ for every proper normal piecewise smooth field $W$ along $\gamma$
if and only if $V$ is a Jacobi field.
\end{corollary}
\section{Geodesics Do Not Minimize Past Conjugate Points}

A geodesic segment $\gamma\colon[a,c]\to M$ \emph{has a conjugate point} if
$\gamma(b)$ is conjugate to $\gamma(a)$ along $\gamma$ for some
$b\in(a,c]$; it has an \emph{interior conjugate point} if one occurs for
$b\in(a,c)$.

\begin{theorem}[Geodesics Do Not Minimize Past Conjugate Points]
\label{thm:no-minimize-past-conjugate}
\dcref{ch10:10.26}
\uses{thm:minimizing-geodesic-unit-speed}
Let $(M,g)$ be Riemannian and let $\gamma$ be a unit-speed geodesic segment
from $p$ to $q$. If $\gamma$ has an interior conjugate point, there is a proper
normal vector field $X$ along $\gamma$ such that $I(X,X)<0$. Consequently,
$\gamma$ is not minimizing.
\end{theorem}

\begin{proof}
\sketch Write $\gamma\colon[a,c]\to M$ and suppose
$a<b<c$ with $\gamma(b)$ conjugate to $\gamma(a)$. Choose a nonzero normal
Jacobi field $J$ with $J(a)=J(b)=0$, and define
\[
V(t)=
\begin{cases}
J(t),&a\leq t\leq b,\\
0,&b\leq t\leq c.
\end{cases}
\]
Then $V$ is proper, normal, and piecewise smooth. Its derivative jump at $b$ is
\[
\Delta D_tV=-D_tJ(b)\ne0,
\]
because $J(b)=D_tJ(b)=0$ would force $J\equiv0$. Choose a smooth proper
normal field $W$ with $W(b)=\Delta D_tV$. Since $V$ is Jacobi on each side of
$b$ and $V(b)=0$, the piecewise index formula gives
\[
I(V,V)=0,
\qquad
I(V,W)=-\langle\Delta D_tV,W(b)\rangle=-|W(b)|^2.
\]
For $X_\varepsilon=V+\varepsilon W$,
\[
I(X_\varepsilon,X_\varepsilon)
=-2\varepsilon|W(b)|^2+\varepsilon^2I(W,W)<0
\]
for all sufficiently small $\varepsilon>0$.
\end{proof}

\begin{lemma}
\label{lem:jacobi-field-constant-pairing}
\dcref{ch10:10.27}
If $J_1,J_2$ are Jacobi fields along a geodesic segment $\gamma$, then
\[
\langle D_tJ_1,J_2\rangle-\langle J_1,D_tJ_2\rangle
\]
is constant along $\gamma$.
\end{lemma}

\begin{proof}
\sketch Differentiate the displayed pairing. The mixed first-derivative terms
cancel, and the Jacobi equation gives
\begin{align*}
\frac d{dt}\bigl(
\langle D_tJ_1,J_2\rangle-\langle J_1,D_tJ_2\rangle
\bigr)
&=-\operatorname{Rm}(J_1,\gamma',\gamma',J_2)\\
&\quad+\operatorname{Rm}(J_2,\gamma',\gamma',J_1)=0
\end{align*}
by the curvature symmetries.
\end{proof}

\begin{theorem}
\label{thm:geodesic-without-conjugate-points}
\dcref{ch10:10.28}
\uses{lem:jacobi-field-constant-pairing}
Let $\gamma\colon[a,b]\to M$ be a unit-speed geodesic segment without
interior conjugate points. Every proper normal piecewise smooth vector field
$V$ along $\gamma$ satisfies $I(V,V)\geq0$, with equality if and only if $V$
is a Jacobi field. If $\gamma(b)$ is not conjugate to $\gamma(a)$, then
$I(V,V)>0$ for every nonzero such $V$.
\end{theorem}

\begin{proof}
\sketch Translate the parameter so that $a=0$, put $p=\gamma(0)$, and choose
an orthonormal basis $(w_1,\ldots,w_n)$ of $T_pM$ with
$w_1=\gamma'(0)$. For $2\leq\alpha\leq n$, let $J_\alpha$ be the normal
Jacobi field satisfying
\[
J_\alpha(0)=0,
\qquad D_tJ_\alpha(0)=w_\alpha.
\]
Absence of interior conjugate points makes $(J_\alpha(t))_{\alpha=2}^n$ a
basis of $\gamma'(t)^\perp$ for $0<t<b$. Thus
\begin{equation}
V(t)=v^\alpha(t)J_\alpha(t), \tag{10.25}
\end{equation}
for piecewise smooth real functions $v^\alpha$. In normal coordinates at $p$,
$J_\alpha(t)=t\,\partial_{x^\alpha}|_{\gamma(t)}$ near $0$; Taylor's theorem,
together with $V(0)=V(b)=0$, gives piecewise smooth endpoint extensions of
the coefficients.

On every smooth subinterval, set
\[
X=v^\alpha D_tJ_\alpha,
\qquad
Y=\dot v^\alpha J_\alpha,
\qquad D_tV=X+Y.
\]
The key identity is
\begin{equation}
|D_tV|^2-\operatorname{Rm}(V,\gamma',\gamma',V)
=\frac d{dt}\langle V,X\rangle+|Y|^2. \tag{10.26}
\end{equation}
Indeed,
\[
\frac d{dt}\langle V,X\rangle
=\langle X+Y,X\rangle+\langle V,D_tX\rangle,
\qquad (10.27)
\]
and the Jacobi equation gives
\[
D_tX=\dot v^\alpha D_tJ_\alpha-R(V,\gamma')\gamma'.
\]
The preceding lemma, evaluated at $t=0$, implies
\[
\langle D_tJ_\alpha,J_\beta\rangle
=\langle J_\alpha,D_tJ_\beta\rangle.
\]
Consequently,
\[
\langle D_tX,V\rangle
=\langle Y,X\rangle-\operatorname{Rm}(V,\gamma',\gamma',V),
\qquad (10.28)
\]
and substitution into (10.27) yields (10.26).

Integrating (10.26) over an admissible partition, the interior boundary terms
cancel because $V$ and $X$ are continuous, and the endpoint terms vanish
because $V(0)=V(b)=0$. Hence
\[
I(V,V)=\int_0^b|Y|^2\,dt\geq0.
\]
Equality holds exactly when each $\dot v^\alpha$ vanishes, equivalently when
$V$ is a constant linear combination of the $J_\alpha$ and hence a Jacobi
field. If the endpoint is not conjugate to $p$, the only proper Jacobi field is
zero, proving positive definiteness on nonzero fields.
\end{proof}

\begin{corollary}
\label{cor:conjugate-point-not-minimizing}
\dcref{ch10:10.29}
\uses{thm:no-minimize-past-conjugate,thm:geodesic-without-conjugate-points}
Let $\gamma\colon[a,b]\to M$ be a unit-speed geodesic segment.
\begin{enumerate}
\item[(a)] If $\gamma$ has an interior conjugate point, it is not minimizing.
\item[(b)] If the endpoints are conjugate but there are no interior conjugate
points, then every proper normal variation $\Gamma$ satisfies
$L(\Gamma_s)>L(\gamma)$ for all sufficiently small nonzero $s$, unless its
variation field is a Jacobi field.
\item[(c)] If $\gamma$ has no conjugate points, the same strict inequality
holds for every proper normal variation with nonzero variation field and all
sufficiently small nonzero $s$.
\end{enumerate}
\end{corollary}

The \emph{index} of a geodesic segment is the largest dimension of a vector
space of proper normal fields on which $I$ is negative definite. The Morse
index theorem states that this index is finite and equals the number of
interior conjugate points counted with multiplicity.

\section{Cut Points}

The flat cylinder supplies a counterexample to the converse of
Theorem~\ref{thm:no-minimize-past-conjugate}: it has no conjugate points, but a
geodesic segment winding more than halfway around the circle factor is not
minimizing.

Let $(M,g)$ be complete and connected, $p\in M$, and $v\in T_pM$. Define the
\emph{cut time}
\[
t_{\mathrm{cut}}(p,v)
=\sup\{b>0:\gamma_v|_{[0,b]}\text{ is minimizing}\}.
\]
It is positive and may equal $\infty$; for $v=0$ it is $\infty$. If
$t_{\mathrm{cut}}(p,v)<\infty$, the \emph{cut point} of $p$ along $\gamma_v$
is
\[
\gamma_v(t_{\mathrm{cut}}(p,v)).
\]
The \emph{cut locus} $\operatorname{Cut}(p)$ is the set of cut points along all
geodesics from $p$. For $\lambda>0$ and $v\ne0$,
\[
t_{\mathrm{cut}}(p,\lambda v)=\lambda^{-1}t_{\mathrm{cut}}(p,v),
\]
so the cut point is unchanged by positive rescaling and it suffices to consider
unit vectors. Any cut point occurs no later than the first conjugate point.

\begin{example}[Cut Loci]
\label{ex:cut-loci}
\dcref{ch10:10.30}
\begin{enumerate}
\item[(a)] For the round sphere $(S^n(R),\bar g_R)$,
$\operatorname{Cut}(p)=\{-p\}$.
\item[(b)] For $S^n(R)\times\mathbb R^m$ with the product metric,
\[
\operatorname{Cut}(p,x)=\{-p\}\times\mathbb R^m.
\]
\end{enumerate}
\end{example}

\begin{exercise}
\label{exer:verify-cut-loci-claims}
\dcref{ch10:10.31}
\uses{ex:cut-loci}
Verify both cut-locus formulas in the preceding example.
\end{exercise}

\begin{proposition}[Properties of Cut Times]
\label{prop:properties-of-cut-times}
\dcref{ch10:10.32}
\uses{thm:no-minimize-past-conjugate}
Let $(M,g)$ be complete and connected, let $p\in M$, let
$v\in T_pM$ be a unit vector, and put
$c=t_{\mathrm{cut}}(p,v)\in(0,\infty]$.
\begin{enumerate}
\item[(a)] If $0<b<c$, then $\gamma_v|_{[0,b]}$ has no conjugate points and
is the unique unit-speed minimizing curve between its endpoints.
\item[(b)] If $c<\infty$, then $\gamma_v|_{[0,c]}$ is minimizing, and at
least one of the following holds:
\begin{itemize}
\item $\gamma_v(c)$ is conjugate to $p$ along $\gamma_v$;
\item at least two unit-speed minimizing geodesics join $p$ to $\gamma_v(c)$.
\end{itemize}
\end{enumerate}

\begin{proof}
\sketch If $0<b<c$, choose $b'$ with $b<b'<c$ such that
$\gamma_v|_{[0,b']}$ is minimizing. An interior conjugate point at or before
$b$ would contradict
Theorem~\ref{thm:no-minimize-past-conjugate}. A shorter curve to
$\gamma_v(b)$ could be concatenated with $\gamma_v|_{[b,b']}$ to shorten the
segment to $\gamma_v(b')$, so $\gamma_v|_{[0,b]}$ is minimizing. If another
unit-speed minimizer reached $\gamma_v(b)$, concatenating it with the remaining
geodesic segment would give a minimizing broken curve with a corner, impossible
for a minimizing curve. This proves (a).

Assume $c<\infty$. Choose $b_i\nearrow c$ with
$\gamma_v|_{[0,b_i]}$ minimizing. Continuity of distance gives
\[
d_g(p,\gamma_v(c))=\lim_i b_i=c,
\]
so the segment through time $c$ is minimizing. Suppose its endpoint is not
conjugate. Choose $b_i\searrow c$. Since the segments through $b_i$ are not
minimizing, let
$\sigma_i\colon[0,a_i]\to M$ be unit-speed minimizers from $p$ to
$\gamma_v(b_i)$ with $a_i<b_i$, and put $w_i=\sigma_i'(0)$. After taking
subsequences,
\[
w_i\to w,
\qquad a_i\to a.
\]
Continuity of $\exp_p$ and distance gives
\[
\exp_p(aw)=\gamma_v(c),
\qquad a=c.
\]
Thus $t\mapsto\exp_p(tw)$ is another minimizing segment to $\gamma_v(c)$.
Because $cv$ is a regular point of $\exp_p$, the map is injective on a
neighborhood of $cv$. The distinct vectors $a_iw_i$ and $b_iv$ have the same
image; since $b_iv\to cv$, the vectors $a_iw_i$ cannot also converge to $cv$.
Hence $cw\ne cv$ and the limiting minimizer is distinct from $\gamma_v$.
\end{proof}
\end{proposition}

Let
\[
UTM=\{(p,v)\in TM:|v|_g=1\}.
\]
The extended interval $(0,\infty]$ carries its usual topology, so continuity at
$\infty$ has the standard infinite-limit meaning.

\begin{theorem}
\label{thm:cut-time-continuous}
\dcref{ch10:10.33}
\uses{prop:critical-points-conjugate}
For a complete, connected Riemannian manifold, the function
\[
t_{\mathrm{cut}}\colon UTM\longrightarrow(0,\infty]
\]
is continuous.

\begin{proof}
\sketch Let $(p_i,v_i)\to(p,v)$ in $UTM$, put
$c_i=t_{\mathrm{cut}}(p_i,v_i)$, and set
\[
b=\liminf_i c_i,
\qquad c=\limsup_i c_i.
\]
We prove
\[
c\leq t_{\mathrm{cut}}(p,v)\leq b.
\]
For the first inequality, pass to a subsequence with $c_i\to c$. If
$c<\infty$, each $\gamma_{v_i}|_{[0,c_i]}$ is minimizing, and continuity of
the exponential map and distance yields
\[
d_g(p,\exp_p(cv))
=\lim_i d_g(p_i,\exp_{p_i}(c_iv_i))
=c.
\]
Thus $\gamma_v$ minimizes through $c$. If $c=\infty$, the same argument on
every fixed interval $[0,c_0]$ shows that $t_{\mathrm{cut}}(p,v)=\infty$.

For the second inequality, assume $b<\infty$ and pass to a subsequence with
$c_i\to b$. Proposition~\ref{prop:properties-of-cut-times} gives infinitely
many indices of one of two types. If $\gamma_{v_i}(c_i)$ is conjugate to $p_i$,
then the determinant characterizing criticality of the fiberwise exponential
map vanishes in the limit. Hence $\gamma_v(b)$ is conjugate to $p$, and
Theorem~\ref{thm:no-minimize-past-conjugate} gives
$t_{\mathrm{cut}}(p,v)\leq b$.

Otherwise choose second unit-speed minimizing geodesics
$\gamma_{w_i}$ from $p_i$ to $\gamma_{v_i}(c_i)$ and pass to a subsequence
with $(p_i,w_i)\to(p,w)$. Both limiting segments reach $\gamma_v(b)$ and are
minimizing. The case $b=0$ is impossible, because then both $c_iv_i$ and
$c_iw_i$ converge to the zero section, where the total exponential map is
locally injective. If the endpoint at time $b$ is conjugate, the preceding
argument applies. If not, $bv$ is a regular point of $\exp_p$. Were $w=v$,
the inverse function
theorem applied to
\[
(x,u)\longmapsto(x,\exp_xu)
\]
near $(p,bv)$ would make the two nearby vectors $c_iv_i$ and $c_iw_i$ with
the same image equal, contradicting the choice of a second geodesic. Thus
$w\ne v$. Part (a) of Proposition~\ref{prop:properties-of-cut-times} then
forces $t_{\mathrm{cut}}(p,v)\leq b$.
\end{proof}
\end{theorem}

For $p\in M$, define
\[
\operatorname{TCL}(p)
=\{t_{\mathrm{cut}}(p,u)u:
u\in T_pM,\ |u|=1,\ t_{\mathrm{cut}}(p,u)<\infty\}
\]
and
\[
\operatorname{ID}(p)
=\{ru:u\in T_pM,\ |u|=1,\ 0<r<t_{\mathrm{cut}}(p,u)\}\cup\{0\}.
\]
These are the \emph{tangent cut locus} and \emph{injectivity domain}. They
satisfy
\[
\operatorname{TCL}(p)=\partial\operatorname{ID}(p),
\qquad
\operatorname{Cut}(p)=\exp_p(\operatorname{TCL}(p)).
\]

\begin{theorem}
\label{thm:cut-locus-properties}
\dcref{ch10:10.34}
\uses{thm:cut-time-continuous}
Let $(M,g)$ be complete and connected, and let $p\in M$.
\begin{enumerate}
\item[(a)] The set $\operatorname{Cut}(p)$ is closed and has measure zero.
\item[(b)] The restriction of $\exp_p$ to
$\overline{\operatorname{ID}(p)}$ is surjective onto $M$.
\item[(c)] The restriction of $\exp_p$ to $\operatorname{ID}(p)$ is a
diffeomorphism onto $M\setminus\operatorname{Cut}(p)$.
\end{enumerate}

\begin{proof}
\sketch If $q_i\in\operatorname{Cut}(p)$ and $q_i\to q$, write
\[
q_i=\exp_p(t_{\mathrm{cut}}(p,v_i)v_i),
\qquad |v_i|=1.
\]
After a subsequence, $v_i\to v$. The cut times equal $d_g(p,q_i)$ and are
bounded, so continuity of cut time and the exponential map gives
\[
q=\exp_p(t_{\mathrm{cut}}(p,v)v)\in\operatorname{Cut}(p).
\]
Thus the cut locus is closed.

In polar coordinates on $T_pM$, the finite part of
$\operatorname{TCL}(p)$ is locally the graph
$r=t_{\mathrm{cut}}(p,\theta)$ of a continuous function. Such graphs have
measure zero; a finite angular chart cover and a countable radial exhaustion
show that $\operatorname{TCL}(p)$ has measure zero. Smooth maps preserve
measure-zero sets locally, so
$\operatorname{Cut}(p)=\exp_p(\operatorname{TCL}(p))$ has measure zero.

Every point of $M$ is reached by a minimizing geodesic from $p$, whose initial
vector lies in $\overline{\operatorname{ID}(p)}$; this proves (b). The
definitions give
\[
\exp_p(\operatorname{ID}(p))=M\setminus\operatorname{Cut}(p).
\]
Before cut time, minimizing geodesics are unique and there are no conjugate
points. Hence $\exp_p$ is injective and a local diffeomorphism on
$\operatorname{ID}(p)$, proving (c).
\end{proof}
\end{theorem}

\begin{corollary}
\label{cor:compact-manifold-quotient}
\dcref{ch10:10.35}
\uses{thm:cut-locus-properties}
Every compact, connected, smooth $n$-manifold is homeomorphic to a quotient of
the closed ball $\overline{\mathbb B}^n$ by an equivalence relation that
identifies points only on its boundary.

\begin{proof}
\sketch Choose a Riemannian metric on $M$ and $p\in M$. Every unit vector in
$T_pM$ has finite cut time bounded by $\operatorname{diam}M$. On the closed
unit ball $\overline B_1(0)\subseteq T_pM$, define
\[
f(v)=
\begin{cases}
t_{\mathrm{cut}}\left(p,\dfrac v{|v|}\right)v,&v\ne0,\\[2mm]
0,&v=0.
\end{cases}
\]
Continuity of cut time, together with its uniform diameter bound near $v=0$,
makes $f$ continuous. It is a bijection
\[
f\colon\overline B_1(0)\longrightarrow
\overline{\operatorname{ID}(p)},
\]
and hence a homeomorphism; it maps the boundary sphere onto
$\operatorname{TCL}(p)$. The exponential map is a continuous surjection from
$\overline{\operatorname{ID}(p)}$ to $M$, so compactness makes it a quotient
map. It is injective on the interior, and interior points and tangent cut points
have disjoint images. Therefore only boundary points are identified. An
orthonormal basis identifies $\overline B_1(0)$ with
$\overline{\mathbb B}^n$.
\end{proof}
\end{corollary}

\begin{proposition}
\label{prop:injectivity-radius-cut-locus}
\dcref{ch10:10.36}
\uses{thm:cut-locus-properties,prop:radial-geodesic-unique-minimizer}
Let $(M,g)$ be complete and connected. For each $p\in M$, the injectivity
radius at $p$ is the distance from $p$ to $\operatorname{Cut}(p)$ when the cut
locus is nonempty, and is infinite otherwise.

\begin{proof}
\sketch Let
\[
d=d_g(p,\operatorname{Cut}(p)),
\]
with $d=\infty$ when the cut locus is empty. For
$a\in(0,\infty]$, let $B_a\subseteq T_pM$ be the open radius-$a$ ball. If
$a\leq d$, then $B_a\subseteq\operatorname{ID}(p)$, so
Theorem~\ref{thm:cut-locus-properties} makes $\exp_p|_{B_a}$ a diffeomorphism
onto its image. If $a>d$, choose a cut point $q$ with $d_g(p,q)<a$. A
diffeomorphic exponential map on $B_a$ would make every radial segment inside
that normal ball uniquely minimizing, contradicting that $q$ is a cut point.
Thus the admissible radii are exactly those not exceeding $d$.
\end{proof}
\end{proposition}

\begin{proposition}
\label{prop:continuity-injectivity-radius}
\dcref{ch10:10.37}
\uses{prop:properties-of-cut-times,thm:cut-time-continuous}
For a complete, connected Riemannian manifold, the function
\[
\operatorname{inj}\colon M\longrightarrow(0,\infty]
\]
is continuous.

\begin{proof}
\sketch For every $p$,
\[
\operatorname{inj}(p)
=\inf\{t_{\mathrm{cut}}(p,v):v\in T_pM,\ |v|=1\}.
\]
The infimum is attained because cut time is continuous on the compact unit
sphere. Let $p_i\to p$ and set
\[
b=\liminf_i\operatorname{inj}(p_i),
\qquad
c=\limsup_i\operatorname{inj}(p_i).
\]
Choose a subsequence with $\operatorname{inj}(p_i)\to b$ and unit minimizers
$v_i$. In a local orthonormal frame, pass to a further subsequence with
$(p_i,v_i)\to(p,v)$. Continuity of cut time gives
\[
\operatorname{inj}(p)
\leq t_{\mathrm{cut}}(p,v)
=b.
\]

For the reverse inequality, take a subsequence with
$\operatorname{inj}(p_i)\to c$. If $\operatorname{inj}(p)<c$, choose
$c_0$ with
\[
\operatorname{inj}(p)<c_0<c
\]
and a unit vector $w\in T_pM$ attaining
$t_{\mathrm{cut}}(p,w)=\operatorname{inj}(p)$. Extend $w$ to unit vectors
$w_i\in T_{p_i}M$ with $(p_i,w_i)\to(p,w)$. Then
\[
t_{\mathrm{cut}}(p_i,w_i)\longrightarrow\operatorname{inj}(p),
\]
so for large $i$,
\[
t_{\mathrm{cut}}(p_i,w_i)<c_0<\operatorname{inj}(p_i),
\]
contradicting the defining infimum. Hence
$c\leq\operatorname{inj}(p)\leq b$, proving continuity.
\end{proof}
\end{proposition}
\section{Problems}

\begin{exercise}
\label{prob:taylor-series-metric-tensor}
\dcref{ch10:10-1}

Let $(M,g)$ be Riemannian and let $p\in M$. Prove that in normal coordinates
centered at $p$,
\[
g_{ij}(x)
=\delta_{ij}-\frac13\sum_{k,l}R_{iklj}(p)x^kx^l+O(|x|^3).
\]
\emph{Hint.} For a radial geodesic
$\gamma(t)=(tv^1,\ldots,tv^n)$ and the Jacobi field
$J(t)=tw^i\partial_i|_{\gamma(t)}$, compute the first four derivatives of
$|J(t)|^2$ at $t=0$ in two ways.
\end{exercise}

\begin{exercise}
\label{prob:area-geodesic-disk}
\dcref{ch10:10-2}
\uses{prob:taylor-series-metric-tensor,prob:associated-quadratic-form}

Let $p$ lie in a Riemannian manifold $(M,g)$, let
$\Pi\subseteq T_pM$ be a $2$-plane, and let
$S_\Pi$ be the surface obtained by applying $\exp_p$ to $\Pi$. For small
$r>0$, let $A(r)$ be the induced area of the geodesic disk of radius $r$
about $p$ in $S_\Pi$. Using Problems~10-1 and~8-21, compute the Taylor
expansion of $A(r)$ through order $r^4$, and express $\sec(\Pi)$ as a limit
involving $\pi r^2-A(r)$.
\end{exercise}

\begin{exercise}
\label{prob:jacobi-fields-constant-curvature}
\dcref{ch10:10-3}
\uses{prop:jacobi-constant-curvature}

Extend Proposition~10.12 by finding a basis for all Jacobi fields along a
geodesic in a constant-curvature manifold, including tangential fields and
fields that do not vanish at $0$.
\end{exercise}

\begin{exercise}[Volume formula for spheres]
\label{prob:volume-formula-spheres}
\dcref{ch10:10-4}

For $n\geq0$, prove
\[
\operatorname{Vol}(S^n(R))=\alpha_nR^n,
\]
where
\[
\alpha_n=
\begin{cases}
\dfrac{2^{2k+1}\pi^kk!}{(2k)!},
& n=2k,\ k\in\mathbb Z,\\
\dfrac{2\pi^{k+1}}{k!},
& n=2k+1,\ k\in\mathbb Z.
\end{cases}
\]
The volume of a compact $0$-manifold is its number of points.
\begin{enumerate}
\item[(a)] Reduce to $R=1$.
\item[(b)] Using Corollary~10.17, prove
\[
\operatorname{Vol}(S^{n+1})=\sigma_n\operatorname{Vol}(S^n),
\qquad
\sigma_n=\int_0^\pi(\sin r)^n\,dr.
\]
\item[(c)] Differentiate $(\sin r)^n\cos r$ to prove
\[
\sigma_{n+1}=\frac{n}{n+1}\sigma_{n-1},
\qquad
\sigma_n\sigma_{n-1}=\frac{2\pi}{n}.
\]
\item[(d)] Complete the proof by induction on $k$.
\end{enumerate}
\end{exercise}

\begin{exercise}[Volume of balls in Euclidean space]
\label{prob:volume-euclidean-balls}
\dcref{ch10:10-5}
\uses{prob:volume-formula-spheres}

For $n\geq1$ and $r>0$, prove that the Euclidean ball satisfies
\[
\operatorname{Vol}(B_r(0))=\frac1n\alpha_{n-1}r^n,
\]
with $\alpha_{n-1}$ as in Problem~10-4.
\end{exercise}

\begin{exercise}[Asymptotic volume of geodesic balls]
\label{prob:asymptotic-volume-geodesic-balls}
\dcref{ch10:10-6}
\uses{prob:taylor-series-metric-tensor,prob:associated-quadratic-form}

Let $p$ lie in a Riemannian $n$-manifold. Using Problems~10-1 and~8-21,
prove that as $r\searrow0$,
\[
\operatorname{Vol}(B_r(p))
=\frac1n\alpha_{n-1}r^n
\left(
1-\frac{S(p)r^2}{6(n+2)}+O(r^3)
\right),
\]
where $S(p)$ is scalar curvature and $\alpha_{n-1}$ is as in Problem~10-4.
\end{exercise}

\begin{exercise}[Nonpositive curvature implies no conjugate points]
\label{prob:nonpositive-curvature-conjugate-points}
\dcref{ch10:10-7}

Prove that a Riemannian manifold with nonpositive sectional curvature has no
conjugate points along any geodesic. \emph{Hint.} For a Jacobi field $J$,
examine derivatives of $|J(t)|^2$.
\end{exercise}

\begin{exercise}[Solvability of the two-point boundary problem]
\label{prob:two-point-boundary-problem}
\dcref{ch10:10-8}
\uses{prop:two-point-boundary-problem}

Prove Proposition~10.21 on solvability of the two-point boundary problem.
\end{exercise}

\begin{exercise}[Normal variation of geodesic between submanifolds]
\label{prob:normal-variation-submanifolds}
\dcref{ch10:10-9}

Let $M_1,M_2$ be embedded submanifolds of a Riemannian manifold $(M,g)$.
Let $\gamma:[a,b]\to M$ be a unit-speed geodesic meeting $M_1$ orthogonally
at $a$ and $M_2$ orthogonally at $b$. Let
$\Gamma:K\times[a,b]\to M$ be a normal variation with
$\Gamma(s,a)\in M_1$ and $\Gamma(s,b)\in M_2$, and let $V$ be its variation
field. Prove
\[
\begin{aligned}
\left.\frac{d^2}{ds^2}\right|_{s=0}L_g(\Gamma_s)
&=\int_a^b
\left(
|D_tV|^2-Rm(V,\gamma',\gamma',V)
\right)\,dt\\
&\quad+\left\langle
\mathrm{II}_2(V(b),V(b)),\gamma'(b)
\right\rangle
-\left\langle
\mathrm{II}_1(V(a),V(a)),\gamma'(a)
\right\rangle,
\end{aligned}
\]
where $\mathrm{II}_i$ is the second fundamental form of $M_i$.
\end{exercise}

\begin{exercise}
\label{prob:frankel-geodesic-intersection}
\dcref{ch10:10-10}
\uses{prob:normal-variation-submanifolds}

Prove Frankel's theorem [Fra61]: if $(M,g)$ is complete, connected, and has
positive sectional curvature, and if compact totally geodesic submanifolds
$M_1,M_2\subseteq M$ satisfy
\[
\dim M_1+\dim M_2\geq\dim M,
\]
then $M_1\cap M_2\neq\varnothing$.
\emph{Hint.} Assuming disjointness, choose a shortest geodesic segment between
$M_1$ and $M_2$ and a parallel field along it tangent to the submanifolds at
the endpoints. Apply Problem~10-9 to obtain a contradiction.
\end{exercise}

\begin{exercise}
\label{prob:jacobi-field-index-form}
\dcref{ch10:10-11}
\uses{cor:jacobi-field-index-form,thm:minimizing-geodesic-unit-speed}

Prove Corollary~10.25: $I(V,W)=0$ for every $W$ if and only if $V$ is a
Jacobi field. \emph{Hint.} Adapt the proof of Theorem~6.4.
\end{exercise}

\begin{exercise}
\label{prob:normal-jacobi-minimizes-index}
\dcref{ch10:10-12}

Let $\gamma:[a,b]\to M$ be a unit-speed geodesic segment with no interior
conjugate points. Let $J$ be a normal Jacobi field, and let $V$ be a piecewise
smooth normal field satisfying $V(a)=J(a)$ and $V(b)=J(b)$.
\begin{enumerate}
\item[(a)] Prove $I(V,V)\geq I(J,J)$.
\item[(b)] If additionally $\gamma(b)$ is not conjugate to $\gamma(a)$ along
$\gamma$, prove that equality holds exactly when $V=J$.
\end{enumerate}
\end{exercise}

\begin{exercise}
\label{prob:killing-vector-jacobi-field}
\dcref{ch10:10-13}
\uses{prob:killing-vector-field,prob:killing-vector-field-properties}

Let $X$ be a Killing vector field on a Riemannian manifold $(M,g)$. Prove that
$X|_\gamma$ is a Jacobi field along every geodesic segment $\gamma$.
\end{exercise}

\begin{exercise}
\label{prob:transverse-jacobi-field}
\dcref{ch10:10-14}

Let $P$ be an embedded submanifold of a Riemannian manifold $M$, and let
$\gamma:I\to M$ be a geodesic meeting $P$ orthogonally at $t=a$. A Jacobi
field $J$ along $\gamma$ is \emph{transverse to $P$} if, on every compact
subinterval containing $a$, it is the variation field of a variation through
geodesics meeting $P$ orthogonally at $a$.
\begin{enumerate}
\item[(a)] Prove that $J(t)=(t-a)\gamma'(t)$ is transverse to $P$.
\item[(b)] Prove that a normal Jacobi field $J$ is transverse to $P$ if and
only if
\[
J(a)\in T_{\gamma(a)}P,
\qquad
D_tJ(a)+W_{\gamma'(a)}(J(a))
\perp T_{\gamma(a)}P,
\]
where $W_{\gamma'(a)}$ is the Weingarten map of $P$ in the direction
$\gamma'(a)$.
\item[(c)] If $\dim M=n$, prove that the transverse Jacobi fields form an
$n$-dimensional subspace of $\mathcal J(\gamma)$ and the transverse normal
Jacobi fields form an $(n-1)$-dimensional subspace.
\end{enumerate}
\end{exercise}

\begin{exercise}[Normal Jacobi fields in semigeodetic coordinates]
\label{prob:jacobi-semigeodetic-transverse}
\dcref{ch10:10-15}
\uses{prop:characterizations-semigeodetic-coordinates}

Let $(x^1,\ldots,x^n)$ be semigeodetic coordinates on an open set $U$ in a
Riemannian $n$-manifold, and let
\[
\gamma(t)=(x^1,\ldots,x^{n-1},t)
\]
be an $x^n$-coordinate curve. For constants
$a^1,\ldots,a^{n-1}$, prove that
\[
J(t)=\sum_{i=1}^{n-1}
a^i\frac{\partial}{\partial x^i}\bigg|_{\gamma(t)}
\]
is a normal Jacobi field transverse to every level set of $x^n$.
\end{exercise}

\begin{exercise}[Jacobi fields via Fermi coordinates]
\label{prob:jacobi-fermi-transverse}
\dcref{ch10:10-16}
\uses{prop:fermi-coordinates-properties}

Let $P$ be an embedded $k$-submanifold of a Riemannian $n$-manifold, and let
$(x^1,\ldots,x^k,v^1,\ldots,v^{n-k})$ be Fermi coordinates on $U_0$. For
fixed $(x^i,v^j)$, set
\[
\gamma(t)=(x^1,\ldots,x^k,tv^1,\ldots,tv^{n-k}).
\]
This is a geodesic meeting $P$ orthogonally at $t=0$. Prove that its Jacobi
fields transverse to $P$ are exactly
\[
J(t)
=\sum_{i=1}^ka^i\frac{\partial}{\partial x^i}\bigg|_{\gamma(t)}
+\sum_{j=1}^{n-k}tb^j
\frac{\partial}{\partial v^j}\bigg|_{\gamma(t)}
\]
for arbitrary constants $a^1,\ldots,a^k,b^1,\ldots,b^{n-k}$.
\end{exercise}

\begin{exercise}[Tangent vectors via transverse Jacobi fields]
\label{prob:tangent-vector-jacobi-normal-neighborhood}
\dcref{ch10:10-17}
\uses{prob:jacobi-fermi-transverse}

Let $U$ be a normal neighborhood of an embedded submanifold $P$ in a
Riemannian manifold. For every $q\in U\setminus P$ and $w\in T_qM$, prove
that there are a geodesic $\gamma:[0,b]\to U$ and a Jacobi field $J$ such
that
\[
\gamma(0)\in P,\qquad
\gamma'(0)\perp T_{\gamma(0)}P,\qquad
\gamma(b)=q,
\]
$J$ is transverse to $P$, and $J(b)=w$.
\emph{Hint.} Use Problem~10-16.
\end{exercise}

\begin{exercise}[Jacobi field in warped product manifold]
\label{prob:jacobi-warped-product}
\dcref{ch10:10-18}
\uses{prob:geodesic-warped-product}

Let $(M_1,g_1)$ and $(M_2,g_2)$ be Riemannian, let
$f:M_1\to\mathbb R^+$ be smooth and positive, and consider
$M_1\times_fM_2$. Fix $q_0\in M_2$, a $g_1$-geodesic
$\gamma:I\to M_1$, and $w\in T_{q_0}M_2$. The curve
\[
\overline\gamma(t)=(\gamma(t),q_0)
\]
is a geodesic orthogonal to every $\{\gamma(t)\}\times M_2$. Prove that
\[
J(t)=(0,w)\in T_{\gamma(t)}M_1\oplus T_{q_0}M_2
\]
is a Jacobi field transverse to every such submanifold.
\end{exercise}

\begin{exercise}
\label{prob:focal-point-characterization}
\dcref{ch10:10-19}
\uses{prob:transverse-jacobi-field}

Let $P$ be embedded in a Riemannian manifold $(M,g)$. A point $q\in M$ is a
\emph{focal point} of $P$ if it is a critical value of the normal exponential
map
\[
E:\mathcal E_P\to M.
\]
Prove that this holds if and only if there are a geodesic
$\gamma:[0,b]\to M$ starting normal to $P$ and ending at $q$, and a nonzero
Jacobi field $J\in\mathcal J(\gamma)$ transverse to $P$ with $J(b)=0$. In
this case, $q$ is a \emph{focal point of $P$ along $\gamma$}.
\end{exercise}

\begin{exercise}
\label{prob:nonpositive-curvature-focal-points}
\dcref{ch10:10-20}
\uses{prob:nonpositive-curvature-conjugate-points}

Let $(M,g)$ have nonpositive sectional curvature and let
$P\subseteq M$ be an embedded totally geodesic submanifold. Prove that $P$
has no focal points. \emph{Hint.} Use Problem~10-7.
\end{exercise}

\begin{exercise}
\label{prob:cut-locus-flat-torus}
\dcref{ch10:10-21}
\uses{ex:n-torus-riemannian-submanifold,ex:n-torus-flat-metric}

Determine the cut locus of an arbitrary point of the flat torus
$\mathbb T^n$ from Examples~2.21 and~7.1.
\end{exercise}

\begin{exercise}
\label{prob:cut-locus-homotopy-equivalence}
\dcref{ch10:10-22}

Let $(M,g)$ be connected and compact, let $p\in M$, and let
$C=\operatorname{Cut}(p)$. Prove that
$C$ is homotopy equivalent to $M\setminus\{p\}$.
\end{exercise}

\begin{exercise}
\label{prob:distance-cut-locus}
\dcref{ch10:10-23}

Let $(M,g)$ be complete, let $p\in M$, and let
$q\in\operatorname{Cut}(p)$ satisfy
\[
d_g(p,q)=d_g\bigl(p,\operatorname{Cut}(p)\bigr).
\]
\begin{enumerate}
\item[(a)] Prove that either $q$ is conjugate to $p$ along a minimizing
geodesic, or exactly two minimizing geodesics
$\gamma_1,\gamma_2:[0,b]\to M$ join $p$ to $q$ and satisfy
\[
\gamma_1'(b)=-\gamma_2'(b).
\]
\item[(b)] Suppose additionally
$\operatorname{inj}(p)=\operatorname{inj}(M)$. If $q$ is not conjugate to
$p$ along any minimizing geodesic, prove that there is a unit-speed closed
geodesic $\gamma:[0,2b]\to M$ such that
\[
\gamma(0)=\gamma(2b)=p,\qquad
\gamma(b)=q,\qquad
b=d_g(p,q).
\]
\end{enumerate}
\end{exercise}

\begin{exercise}
\label{prob:injectivity-radius-exponential}
\dcref{ch10:10-24}

Let $(M,g)$ be complete and let $p\in M$. Prove that
$\operatorname{inj}(p)$ is the radius of the largest open ball in $T_pM$ on
which $\exp_p$ is injective.
\end{exercise}
